\section{The denominator with the norm-choice family hypotheses}

% BEGIN V84A1
\paragraph{V84A1. The information retained by the alternate hypotheses.}
Fix a member $D$, its norm choice $n$, and its
$\mathrm{NormChoiceFacts}(D,n)$. At every localization $i$ and for
every good integral family $R$, the following data are available.
\begin{enumerate}
\item The coefficients of $\det(1-zR_x(g))$ belong to the Hecke
algebra as functions of $x$, the residual characteristic polynomial
at $g$ is independent of $x$, and $R$ is unramified outside the
places above $p=\ell_i$ and the two sets $S_i,Q_i$. After one finite
coefficient extension, every irreducible subrepresentation at every
$x$ is absolutely irreducible and disjoint irreducible
subrepresentations at the same $x$ are inequivalent.
\item Off the two type primes, every $R_x$ is crystalline above $p$,
with the specified level-zero jumps. At either type prime every
$R_x$ is de Rham there with the same labelled jump formula.
At parameters satisfying $\mathrm{AdmPar}$ one has the sign
dictionary, the four multiplicity readings and the
crystalline--Weil--Deligne tie of $\mathrm{AdmJAlt}$.
\item The quadratic inertia identities and the Hecke polynomial at
a type place in $S_i$, the off-support criterion for a global
$\chi_1$-line, and the good-row image and rank identities have
the original $\mathrm{Conclusion2}$ quantifiers.
\end{enumerate}
Indeed $\mathrm{provisionFortyOne}$ supplies $G1,G2a,G3,G4$ in the
alternate definitions. Take the second factor of $G3$ for the
residual polynomial. Its first factor requires a guard, but the
second factor of $\mathrm{DeRhamJumpsDistinctCondition}$ gives
the entire characteristic-zero distinctness assertion for the
chosen good family without that guard. The first factor of the
latter condition gives exactly the type-prime reading. $G2a$ is
the off-type crystalline and jump assertion itself. Take the
remaining reading fields from $\mathrm{admJ}(L)$ at the supplied
parameter. For (3) project to $\mathrm{ConclusionG5}$,
$\mathrm{ConclusionG6}$ and $\mathrm{ConclusionZ}$, respectively.

For clarity about the comparison, the seven fields extending the
core to norm-choice facts (localization uniqueness, large-prime
conditions, coefficient duties, norm-choice tie, good triples, same
twist and goodness of all the specified signatures) are unchanged.
In the core, $\mathrm{Conclusion2}$, residual irreducibility of the
rank block, weight and character dictionaries, unramifiedness of
the twist, depth-zero condition, spherical base change and
$\mathrm{LamUsed}$ are unchanged. The changed entries are the
row record, admissibility, first conclusion, residual tame labels,
de Rham conditions and $\mathrm{admJ}$. Inside the row record
the changes are in its first part, concerning Galois attachment,
the cuspidal bound and the former all-place invariant condition.
The alternate attachment gives a relation of local parameters
away from $p$, giving equality whenever the automorphic parameter
is unramified. Its clause above $p$ has a rank-one or local-parameter
guard. The alternate cuspidal bound gives temperedness in rank one
and $q^{-1/2}<|\alpha|<q^{1/2}$ on unramified roots. The former
all-place invariant condition is replaced by the true predicate.
The alternate residual tame label assertion has a squarefreeness
premise at all the type labels. None of (1)--(3) requires removing
that premise. Finally $\mathrm{ExtraAlt}$ keeps the same
commuting-endomorphism field. Its paired-character inequality
has the same right side and omits the argument $\mathrm{SetSQ}$.
It therefore implies the ordinary paired inequality by applying
it to the other arguments. These projections specify the inputs
used below without an implication between the complete core records.
% END V84A1

% BEGIN V84A2
\paragraph{V84A2. Constructing the alternate setup on its own member.}
Assume $\mathrm{Hypotheses}_2$, $\mathrm{TwGood}(\mathrm{Fm})$ and
$\mathrm{FactsFAlt}(\mathrm{Fm},C,C_T)$. For each real $C_2$ there is
a cutoff such that for every subsequent $x>\mathrm{Fm}.x_1$, each
member at $x$ satisfying $\mathrm{ShaFinOdd}$, and every prescribed
finite set $A$, there are $H,F,\mathrm{Cmp}$ and a hub $X$ on this
member with
\[
 X.\mathrm{exc}=A,\qquad \mathrm{AdmJAlt}(L),\qquad
 \mathrm{SetupAlt}(L,n,H,X,C_2).
\]

To construct the point representations, fix a regular algebraic
conjugate-self-dual cuspidal part of rank $a\ge1$. The alternate
attachment supplies an actual continuous semisimple representation.
Discard the finitely many places over its coefficient prime and
the finite set of ramified automorphic places. At each remaining
place its normalized parameter is unramified because the
normalizing twist is unramified. The equality clause of the
local relation now identifies this parameter with the attached
Weil--Deligne parameter. It too is unramified. By the tied
Frobenius polynomial, the trace of this representation at an
arithmetic Frobenius equals the sum of the Satake parameters
times $(\mathrm Nw)^{(1-a)/2}$ under the specified coefficient
identification. This proves precisely the cofinite trace input
DFD4 supplies in the ordinary construction.

Choose the good representation family promised at each index by
the first existence clause of $\mathrm{provisionFortyOne}$.
Using these trace representations, the two parity cases of DFD5
give all shifted blocks. DFD3 and DFD6 identify their traces on
the retained Hecke point and produce its character splitting.
The rational descent and saturation in DFD1--DFD2 give a
continuous complement over this very integral coefficient ring.
DFD7--DFD8 assemble the point family and its canonical hub for
the given twist, using zero as error function and $A$ as the
constant exceptional-set function. This is the construction of
DFD9 with the trace input proved above.

Use the three cutoffs of F78D2 for the family parameter, small
primes at constant one and the good twist. Choose $H,F,\mathrm{Cmp}$
by L78W13 before forming that hub. The good twist gives its twist
index. The first conclusion about theta lifts and its good-row
image, used by DFD10, are respectively $\mathrm{Conclusion42}$
in the alternate first conclusion and $\mathrm{ConclusionZ}$ in
the second. For the actual point representations the norm-choice
tie, good triples, same twist and residual irreducibility are the
fields in V84A1. Thus SSC15--SSC16 and DFD11 give killing on
$X.\rho$. The family size cutoff gives $8\le L.\mathrm{sc}.L$ and
$60\le D.\mathrm{host.splitBound}$. DFD13 gives $\mathrm{AdmPar}$
and the polynomial parameter bounds, hence $\mathrm{AdmJAlt}$ by
the member's norm-choice facts. These data and the supplied
finiteness fill every field of $\mathrm{SetupAlt}$.
% END V84A2

% BEGIN V84A3
\paragraph{V84A3. The jump, global Hecke and type-prime tests.}
Fix any alternate setup, its actual $R=X.r$, $\rho=X.\rho$ and
the native member parameters. The following verifications also
hold for arbitrary good $R$ and specified rank-$r$ complements
when these are quantified in the corresponding test.

First, S80J1--S80J2 give all seven jump conditions. Their
calculation uses the label and the sign in
$\mathrm{AdmDict}$ and $\mathrm{AdmRead}$ of V84A1. At a non-type
prime, the crystalline comparison gives the required de Rham
reading. At a type prime, use the supplied de Rham clause.
Subtract the same one-dimensional character weight, then subtract
$\{0,\ldots,t-1\}$ at the split rank block. In both orientations
the resulting rank list and its dual $z\mapsto-1-z$ omit zero.
At a non-type prime, the full family is crystalline, and the
character and rank block are crystalline direct summands at the
chosen split point. The same holds after twisting and coefficient
extension by S80J3. Any rational invariant in the rank block or
its Tate dual would supply a trivial de Rham subobject and hence
a jump of weight zero by strictness. This proves both required
local vanishings above $p$ on this data.

At an index with good support the Hecke coefficient and residual
polynomial conditions of V84A1 and the identical $\mathrm{rhoBar}$
field give the full integral matrix span of D79G1. The coefficient
duties supply the stated decomposition over the coefficient field
in S80J3. At a point outside $Z_i$, a global fixed vector after
dividing by $\chi_1$ would supply the $\chi_1$-line tested by
$\mathrm{ConclusionG6}$. Outside $S_i,Q_i,p$, $G4$ gives the
unramified action, including on the character summand at the
chosen point. In $Q_i$ and at the auxiliary prime in $S_i$,
$\mathrm{ConclusionG5}$ gives the quadratic inertia identity;
at a type place in $S_i$ it gives the annihilating Hecke
polynomial and its unit value. The tame nonidentity and auxiliary
order-$15$ unit calculation of D79G7 depend only on the specified
characters. Thus D79G1--D79G7 give the same determinant, local
and off-support inputs in their respective ranges.

For the type-prime estimate assume also the paired-character
part of $\mathrm{ExtraAlt}$, $0\le C$, $C_T\le b$ and the late
family size bounds. At $p=g_k$, take $s=\lfloor3b\rfloor+2$ and
$S=0$ as in T79W1. The log-ratio and size inequalities in Y79T1
and T79W1 depend only on the two type primes and the member
sizes. At the completely split prime, T79W2 gives the opposite
signs at conjugate places and the sets of jumps for every $R_x$
by the type-prime reading just proved. These are
$\mathrm{SignOK}$ and $\mathrm{HostJ}$. Every $\pi\in Z_i$ has
its canonical semisimple splitting. Killing gives the finite
Selmer predicate at every good point and $\mathrm{ShaFinOdd}$
gives finiteness at $0<e<t$. Substitution in the alternate paired
inequality yields exactly the bound for $\eta_{\rm loc}$ in Y79T1.
In the empty good-support case both the good eigenspace index
and the sum vanish as in that proof. Y79T2--Y79T4 then give
$\mathrm{StarH}$ with the original multiplicities and type
constant $2000(C+1)+5C_1$, where $C_1$ is the uniform parameter
bound. Here the multiplicity and rank bounds used in the summation
follow from $\mathrm{Conclusion42}$, $\mathrm{ConclusionZ}$,
$\mathrm{adm8}$, localization uniqueness and the rank field by LN6
and J78H1. The discarded all-place condition does not occur in
this substitution.
% END V84A3

% BEGIN V84A4
\paragraph{V84A4. The small-prime family bound for every alternate setup.}
Under $\mathrm{Hypotheses}_2$, every family satisfies
$\mathrm{BoundAtTwoThreeFamilyFonAlt}$. In fact a single pair
$(C_s,x_0)$ bounds all its good-point rows at $p\mid Q_S$ or
$p=3$ whenever the quantified alternate setup is supplied.

Let $C_1\ge1$ be the parameter constant of A78P3 for this family.
Take
\[
 C_s\ge18\cdot6000 C_1+2\cdot30000(3C_1+9)+5C_1+1 .
\]
Fix the cutoff for these parameter and size inequalities and
$\log x\ge1$ before any member, functor, hub or $C_2$.
Now take arbitrary $H,F,\mathrm{Cmp},X,C_2$ and a witness
of $\mathrm{SetupAlt}$. At a twisted index with a good point the
setup gives a genuine point family, its killing, the norm-choice
facts, the splitting bound and the finiteness at $0<e<t$ for
every prime, including $2$. Specializing the norm-choice facts
at its admissible parameters supplies the reading data of V84A3.
The argument for this assertion uses neither
$\mathrm{ExtraAlt}$ nor a good twist of the whole family.

For these primes G82F1 gives the split declaration, the levels
$s=\max(1,v_p(Q_S))$ and the cyclotomic inertia element. It uses
only the host and splitting bound. Each prime is outside the two
type primes. V84A3 therefore supplies the full ambient and dual
crystalline families on each tested coefficient extension, their
simple labelled jumps and the rank-block invariant vanishing.
In G82F2 construct the integral matrix span, projectors and
auxiliary inertia relations from the Hecke and off-support tests
of V84A3. Their root differences at the chosen lift have valuation
$\sum_{1\le j<t}(s+v_p(j))$. The small-prime calculation uses
the declaration just chosen and the actual $R,\rho$ in this
setup. The row constant $C_2$, hub exceptional set and stored
middle-range lengths do not enter it.

More precisely, L82K5 supplies the integral lifts for all
surjections onto the rank and character columns and for the
dual row maps, since the crystalline subquotients have the
inherited height intervals. T82K6 supplies the torsion and
invariant estimates on those same representations. One may
use the compatible comparison of L78W13 on $D$ for these local
calculations: the maps and finite local groups in G82F2 are
defined on the Galois lattices of $X$, independently of the
presentation used to calculate them. The determinant argument
and descent in G82F2 now give $\mathrm{Bd}$ with its displayed
error $\mathcal E$. By G82F3,
\[
 \mathcal E\le6000 f m h_B s+
    30000 f^2t^2 n_V\bigl(s+\log_p(2h_B)+1\bigr).
\]
G82F4 bounds this by $\mathrm{midF}(L,C_s)$ with precisely
the choice above. At $p=2$ the level is at least two because
$4\mid Q_S$; at $p=3$ outside $Q_S$ the setup's splitting bound
suffices. These include exactly the two alternatives in
$\mathrm{BoundAtTwoThree}$. Keeping also the other primes of
$Q_S$ yields $\mathrm{StarF}$ on this $X.\rho$ by G82F4's
auxiliary table of zero place lengths and non-place entry
$\mathrm{midF}(L,C_s)$.
% END V84A4

% BEGIN V84A5
\paragraph{V84A5. Fifteen simultaneous fields and the alternate-family denominator.}
Under $\mathrm{Hypotheses}_2$, suppose the given family has
$\mathrm{TwGood}$ and $\mathrm{FactsFAlt}(\mathrm{Fm},C,C_T)$
with $C\ge0$. Then it satisfies $\mathrm{ConclF}$ and
$\mathrm{LinkN}$ with their original late-member quantifiers.

Set $C_2=122$. Choose the common table constant as in A83G2,
including the lower bounds in V84A3--V84A4 and $C_T$. At this
constant choose the cutoff meeting all the size bounds
(including $C_T\le b$), the local and global row bounds and the
cutoff of V84A2. At any subsequent member satisfying
$\mathrm{ShaFinOdd}$, prescribe its arithmetic exceptional set in
V84A2. Keep that very hub in the following verifications.
\begin{itemize}
\item The canonical table, exceptional-set equality, admissible and
polynomial parameters give $\mathrm{num},\mathrm{exc},\mathrm{par}$
by F78D3, X78A3 and A78P3. Positivity is the choice of the
constant and $\mathrm{small}$ is DFD14 at this constant.
$\mathrm{pass}$ follows from LN5 with the given commuting
endomorphisms. Multiplicity and rank were checked in V84A3 and
killing is part of V84A2.
\item For $\mathrm{bb}$, the hypotheses on the twisted rank
lattice at the declared columns in S79P1--S79P5 are its
crystallinity, jump bounds and absence of zero from V84A3.
J79H1's exponent $120$ thus gives the coarse bound $7200$ on
this canonical hub by J79H2. The type field is V84A3.
\item At a large-prime test, fix any of its good families,
points and complements. The full span, auxiliary inertia and
off-support conditions of V84A3 give D79G8's global inputs.
The guard excludes the type primes. The jump lists, crystalline
condition and residual and determinant data in V84A3 are exactly
the member inputs to S80J3--S80J6. J80H1--J80H2 therefore give
the three local cohomology assertions and the congruence estimate
for this arbitrary complement. With the unchanged weight dictionary
and the hub's exceptional set, G79E2 gives $\mathrm{large}$.
\item The primes $q,d$ and the guarded middle range are also
outside the type primes. For $q,d$ apply J81H2 on the rank
and character cuts formed with V84A3's global data. The rank
lifts and character bounds there depend on the crystalline
jump lists and the local fields; both are those already verified.
For a middle-range prime use the same Hecke and inertia data for
the common columns of G82C5 or the empty columns of G82I1.
At its places the cut lattice, labelled weights and finite
cohomology are supplied by the canonical table and V84A3.
The split estimate G82C5 and inert estimates B82R3 and A83G1
then give the actual $\mathrm{StarAD}$ on this hub. Their
constants are the prescribed row constants at $C_2=122$.
This gives $\mathrm{ram}$ and $\mathrm{mid}$. The small-prime
table of V84A4 gives $\mathrm{two}$ for the very same $X.\rho$.
\end{itemize}
This constructs every field of $\mathrm{Concl}$ for the chosen
$X$ within the uniform cutoff, hence $\mathrm{ConclF}$.
Apply the table export of A83G3 to these fields. For the same
member and each allowed lift, under summability of its
$\mathrm{vvTerm}$, the weighted $\eta$-sum satisfies
$\mathrm{LinkN\_at}$ with constant $2000(|C_s|+1)^3$.
Both constants and any enlargement of the export cutoff precede
the member. This proves $\mathrm{LinkN}$ on the supplied alternate
family, including when that family is the one obtained by cutting.
% END V84A5
